\scanpage{142}

% L1–3, continuation
ramené à prouver que ce dernier est connexe. Or il est fibré sur $\underline\Delta_{\{n-1\}}$ (qui est connexe), il suffit donc de prouver

% L4–5, lemma
\uline{Lemme 6} $\underline\Delta_{\{n-1,n\}}^\omega\xrightarrow{\sim}\underline\Delta_{\{n-1\}}$ par la projection canonique, en d'autres termes

% L5–7, left edge
% slanted
\margin{\unsure{indépendt|indépend|indépendant} de tte question de car.}

% L6–8, statement
$\forall F'\in\Delta_{\{n-1\}}(E,q)$, $\exists$ ex.ct$^{\mathrm{mt}}$ un $F\in\Delta_{\{n\}}^\omega(E,q)$ tel que $F\subset F'$. [Ou encore \deleted{\unsure{L}}la flèche de $(\underline\Delta_{\deleted{\unsure{H}}\,\{n-1,n\}})\mid F$ dans $\underline{\mathrm{Arf}}\deleted{\unsure{(E)}}$ est un iso]

% L9–14, proof
\uline{Dém} Par lemme 3, on voit que l'image inverse de $F'$ dans $\underline\Delta_{\{n,n-1\}}$ est $\unsure{\simeq|\text{à}}\Delta_1(\underbrace{F'^\perp/F'}_{\mathrm{rg}\,2})$ donc est un rev. étale de degré 2 de $S$. Il suffit de voir que le morphisme \inserted{de celui-ci} dans $\underline{\mathrm{Arf}}$ commute aux \unsure{autom.|actions} can. de ces rev., ce qui est immédiat.

% L15–20, lower text
Une fois traité le cas de car $\neq2$, \inserted{pour déterminer le nb de comp. \deleted{\unsure{de}} connexes de $\underline\Delta_{\{n\}}$,} on obtient le cas général par un argument de spécialisation (fact. de Stein....). En $=$ temps, \unsure{cela|on} définit $\underline{\mathrm{Arf}}(E,q)$ dans le cas général, par factorisation de Stein de $\underline\Delta_{\{n\}}$, et on

% L21, bottom continuation
voit que cela marche pour tous les $\underset{\substack{H\ni n\,\deleted{\unsure{\{n\}}}}}{\underline\Delta_H}$ \unsure{ou|où|avec}

% bottom left, pencil
% archive sheet number: 141
